%=============================================================
\chapter{The Harmonic Drive Sub-Model}
\label{chap:harmonic}
%=============================================================

\section{Motivation}

Real robotic joints are not directly driven: a gearbox with a high reduction ratio sits between the motor shaft and the output link.
In the UR5, \emph{harmonic-drive} (strain-wave) gearboxes are used.
These components exhibit several non-ideal characteristics that significantly affect the closed-loop performance:
\begin{enumerate}
  \item High gear ratio with near-zero backlash.
  \item Non-linear torsional stiffness (the flex-spline deforms elastically).
  \item Viscous friction at the input shaft.
  \item Asymmetric Coulomb friction (static friction differs for positive and negative rotation).
  \item Motor rotor inertia reflected through the gear ratio.
\end{enumerate}

The harmonic drive datasheets (HFUS-14 and HFUS-20 series, see \cref{tab:hd14,tab:hd20}) supply the rated and peak torques, maximum speeds, and moment-of-inertia values used to set the model parameters.

\begin{table}[H]
  \centering
  \caption{HFUS-14 harmonic drive specifications (selected entries).}
  \label{tab:hd14}
  \begin{tabular}{@{}lcc@{}}
    \toprule
    Parameter & Ratio 50 & Ratio 100 \\
    \midrule
    Repeated peak torque $T_p$ [Nm] & 18 & 28 \\
    Average torque $T_a$ [Nm]       & 6.9 & 11 \\
    Rated torque $T_r$ [Nm]         & 5.4 & 7.8 \\
    Moment of inertia $J$ [$10^{-4}$ kg\,m$^2$] & \multicolumn{2}{c}{0.033--0.091} \\
    \bottomrule
  \end{tabular}
\end{table}

\begin{table}[H]
  \centering
  \caption{HFUS-20 harmonic drive specifications (selected entries).}
  \label{tab:hd20}
  \begin{tabular}{@{}lcc@{}}
    \toprule
    Parameter & Ratio 80 & Ratio 100 \\
    \midrule
    Repeated peak torque $T_p$ [Nm] & 74 & 82 \\
    Average torque $T_a$ [Nm]       & 47 & 49 \\
    Rated torque $T_r$ [Nm]         & 34 & 40 \\
    Moment of inertia $J$ [$10^{-4}$ kg\,m$^2$] & \multicolumn{2}{c}{0.193--0.404} \\
    \bottomrule
  \end{tabular}
\end{table}

\section{Ports and Parameters}


\section{Bond-Graph Structure}

The \texttt{HarmonicDrive} sub-model has:
\begin{itemize}
  \item \texttt{motor}: scalar power port input (motor torque + angular velocity).
  \item \texttt{harmonic\_drive}: scalar power port output (joint torque + output shaft velocity).
\end{itemize}

The topology is shown schematically in the bond graph below.
The key elements are:

\paragraph{Motor inertia (I element).}
The rotor inertia $J_{\text{rotor}}$ (default $8\times10^{-6}$~kg\,m$^2$) stores kinetic energy of the motor shaft:
\begin{equation}
  J\,\dot{\omega}_m = \tau_m - \tau_{\text{friction}} - \tau_{\text{elastic}}/N,
\end{equation}
where $\omega_m$ is the motor angular velocity and $N$ is the gear ratio (stored as negative to account for sign convention, default $N=-0.01$, i.e.\ ratio of 100).

\paragraph{Gear ratio (TF element).}
The transformer element with ratio $N$ enforces:
\begin{equation}
  \tau_m = N \, \tau_{\text{out}}, \qquad \omega_{\text{out}} = N \, \omega_m.
\end{equation}

\paragraph{Torsional stiffness (C element -- non-linear).}
The flex-spline introduces a piecewise-linear torsional compliance parameterised by three stiffness zones:
\begin{equation}
  \tau_s(\varphi) =
  \begin{cases}
    k_1\,\varphi & \text{if } \varphi \leq \varphi_1 = T_1/k_1, \\
    k_2(\varphi-\varphi_1) + T_1 & \text{if } \varphi_1 < \varphi < \varphi_2, \\
    k_3(\varphi-\varphi_2) + T_2 & \text{if } \varphi \geq \varphi_2,
  \end{cases}
\end{equation}
where $\varphi = \int(\omega_{\text{out}} - \omega_{\text{load}})\,dt$ is the elastic twist of the flex-spline, and $k_1 = 4700$, $k_2 = 6100$, $k_3 = 7100$~Nm/rad, $T_1 = 2$~Nm, $T_2 = 6.9$~Nm.

\paragraph{Viscous friction at motor shaft (R element).}
\begin{equation}
  \tau_{\text{visc}} = b_w\,\omega_m, \qquad b_w = 6.4\times10^{-4}\;\text{Nm\,s/rad}.
\end{equation}

\paragraph{Asymmetric Coulomb friction at output shaft (R element).}
Friction at the output shaft is modelled as a smooth approximation to the sign function:
\begin{equation}
  \tau_f(\omega) = \frac{b_{f,+} - b_{f,-}}{2} \cdot \frac{2}{\pi}\arctan(\omega) + (b_{f,+}+b_{f,-}),
\end{equation}
with $b_{f,+} = 0.008$~Nm and $b_{f,-} = -0.007$~Nm.
This formulation avoids numerical stiffness introduced by a discontinuous sign function while faithfully capturing the asymmetry.

\subsection{Causality Assignment}

The \texttt{motor} port is fixed \emph{effort-in}: the motor torque source imposes effort.
The \texttt{harmonic\_drive} port is fixed \emph{effort-out}: the transmission output imposes torque on the joint.
Within the sub-model:
\begin{itemize}
  \item The 1-junction at the motor shaft has the I-element as the \emph{effort-out} element (integral causality), and the TF and R elements as effort-in elements.
  \item The 0-junction at the output shaft is the sum point for the elastic and Coulomb friction torques.
  \item The C element is in integral causality: $\dot{\varphi} = f_C$, $e_C = \tau_s(\varphi)$.
\end{itemize}
The non-linear C element breaks the direct algebraic loop that would exist if a linear C were used in derivative causality.

\begin{placeholder}[FIGURE PLACEHOLDER -- Harmonic Drive Bond-Graph Diagram]
  \centering
  \vspace{3cm}
  \textit{Insert here the bond-graph diagram of the \texttt{HarmonicDrive} sub-model.
  Show the 1-junction at the motor shaft (connecting the motor source, the I-element for rotor inertia, the R-element for viscous damping, and the TF for gear ratio), the 0-junction at the output shaft (connecting the elastic C-element, the asymmetric friction R1-element, and the harmonic\_drive output port), and the inner 1-junction coupling the two.
  Causality strokes should be clearly indicated.}
  \vspace{3cm}
\end{placeholder}